\section{From combinatorics to probability models}

The role of combinatorics in this course is not to memorize more formulas. Its role is to build sample spaces and events precisely.

The general workflow is
\[
\text{situation}
\longrightarrow
\text{recording rule}
\longrightarrow
\text{elementary outcome}
\longrightarrow
\text{sample space}
\longrightarrow
\text{event}
\longrightarrow
\text{probability}.
\]

In later chapters, many probability distributions will arise by recording only a numerical summary of a richer experiment. For example, in \(n\) Bernoulli trials the full elementary outcome may be a sequence such as
\[
(S,F,S,S,F),
\]
while the random variable records only the number of successes:
\[
X=3.
\]
The binomial coefficient appears because many different sequences have the same number of successes.

Thus combinatorics is the bridge between concrete experiments and probability laws. It explains how many elementary outcomes produce the same event, the same count, or the same value of a random variable.

\section{Final checklist}

When solving a counting problem, write down the following before doing arithmetic:

\begin{enumerate}
    \item What is the physical situation?
    \item What information is recorded?
    \item What is one elementary outcome?
    \item What is the sample space?
    \item Are outcomes sequences, sets, arrangements, or count vectors?
    \item Are the elementary outcomes equally likely if probability is being computed?
    \item Which construction or overcounting argument gives the count?
    \item What does the final number mean in the original situation?
\end{enumerate}

This checklist is the main lesson of the chapter. The formulas are useful, but they are useful only when the underlying structure has been discovered first.
